%! Function Model Construction and Application — Unit 1, Lesson 1.14 — Group Activity (Answer Key)
\documentclass[10pt]{article}
\usepackage{precalculus-article}
\usepackage{precalculus-key}   % pulls in -boxes; do NOT also load -boxes
\newcommand{\TallMath}[1]{$\displaystyle #1\rule[-1.4em]{0pt}{3.2em}$}

\begin{document}

\pageheader{Unit 1, Lesson 1.14}{Group Activity --- Answer Key}
\namepartnerperiod

\vspace{0.08in}

\begin{scenariobox}[The Engineering Design Challenge]{sky}
\small
Engineers use function models to make decisions. In each scenario below, you will
\textbf{construct a model} and \textbf{apply it} to answer design questions.
Work through your assigned tier. If you finish early, try the next tier.
\end{scenariobox}

\vspace{0.08in}

\begin{tcolorbox}[colback=white, colframe=black!40,
    title=\textbf{Tier R --- Remediate},
    fonttitle=\bfseries, arc=2mm, left=3mm, right=3mm, top=2mm, bottom=2mm]

\textbf{Scenario:} A rectangular solar panel array has one side against a wall.
The available border is 60 m. The area function has already been set up for you:
\[
  A(w) = -2w^2 + 60w, \quad 0 < w < 30
\]
where $w$ is the width in meters and $A$ is the area in square meters.

\begin{enumerate}[itemsep=10pt]
  \item Evaluate the model at each width. Show your substitution.

        \medskip
        \renewcommand{\arraystretch}{1.6}
        \begin{tabular}{|c|p{5cm}|c|}
        \hline
        $w$ (m) & Work & $A(w)$ (m$^2$) \\
        \hline
        5  & $-2(25)+300$ & \ans{250} \\
        \hline
        10 & $-2(100)+600$ & \ans{400} \\
        \hline
        15 & $-2(225)+900$ & \ans{450} \\
        \hline
        25 & $-2(625)+1500$ & \ans{250} \\
        \hline
        \end{tabular}

  \item Compute the average rate of change of $A$ on the interval $[5, 15]$.

        \vspace{0.06in}
        $\text{ARC} = \dfrac{A(15) - A(5)}{15 - 5} = \dfrac{{\color{keyred}\mathbf{200}}}{{\color{keyred}\mathbf{10}}} = {}$ \ans{20}

        \vspace{0.04in}
        Units: \ans{m$^2$ per m (square meters per meter of width)}

  \item Based on your table in part (1), between which two widths does the maximum area
        appear to occur?

        \vspace{0.06in}
        Between $w = $ \ans{10} m and $w = $ \ans{25} m \ans{(vertex at $w=15$; maximum is 450 m$^2$)}.
\end{enumerate}
\end{tcolorbox}

\vspace{0.08in}

\begin{tcolorbox}[colback=white, colframe=black!40,
    title=\textbf{Tier A --- Apply},
    fonttitle=\bfseries, arc=2mm, left=3mm, right=3mm, top=2mm, bottom=2mm]

\textbf{Scenario:} A storage company charges based on distance from a central depot.
Testing shows that the shipping cost $C$ (in dollars) is inversely proportional to the
square of the distance $d$ (in km) from the depot, with proportionality constant $k = 2000$.

\begin{enumerate}[itemsep=10pt]
  \item Write the rational function model $C(d)$.

        \vspace{0.06in}
        $C(d) = $ \ans{$\dfrac{2000}{d^2}$}

  \item State the domain restriction and explain it in context.

        \vspace{0.06in}
        Domain: \ans{$d > 0$} \quad because \ansline{distance cannot be zero or negative --- the depot must exist at a positive distance from the delivery point.}

  \item Complete the table and describe the trend.

        \medskip
        \renewcommand{\arraystretch}{1.5}
        \begin{tabular}{|c|c|c|c|c|}
        \hline
        $d$ (km) & 2 & 4 & 10 & 20 \\
        \hline
        $C(d)$ (\$) & \ans{500} & \ans{125} & \ans{20} & \ans{5} \\
        \hline
        \end{tabular}

        \vspace{0.06in}
        Trend: As distance increases, cost \ans{decreases (approaches 0)}.

  \item Compute the average rate of change of $C$ on $[2, 4]$. Include units.

        \vspace{0.06in}
        $\text{ARC} = \dfrac{C(4) - C(2)}{4 - 2} = \dfrac{{\color{keyred}\mathbf{-375}}}{{\color{keyred}\mathbf{2}}} = {}$ \ans{$-187.50$}

        Units: \ans{dollars per km (\$/km)}
\end{enumerate}

\begin{teachernote}
ARC $= -187.50$ \$/km: the negative sign indicates cost decreases as distance increases.
Students should note this is the average decrease, not a steady rate (the model is not linear).
\end{teachernote}
\end{tcolorbox}

\vspace{0.08in}

\begin{tcolorbox}[colback=white, colframe=black!40,
    title=\textbf{Tier E --- Extend},
    fonttitle=\bfseries, arc=2mm, left=3mm, right=3mm, top=2mm, bottom=2mm]

\textbf{Scenario:} An engineer collects three data points for a project's profit $P$
(in thousands of dollars) versus number of units produced $x$ (in hundreds):

\medskip
\renewcommand{\arraystretch}{1.3}
\begin{tabular}{|c|c|c|c|}
\hline
$x$ & 1 & 2 & 3 \\
\hline
$P(x)$ & $-1$ & 5 & 5 \\
\hline
\end{tabular}

\begin{enumerate}[itemsep=10pt]
  \item Assume $P(x) = ax^2 + bx + c$. Substitute each data point to write a system
        of three equations in $a$, $b$, $c$.

        \vspace{0.06in}
        $(1,-1)$: \ans{$a + b + c = -1$}

        \vspace{0.04in}
        $(2, 5)$: \ans{$4a + 2b + c = 5$}

        \vspace{0.04in}
        $(3, 5)$: \ans{$9a + 3b + c = 5$}

  \item Solve the system. Show your elimination or substitution steps.

        \vspace{0.08in}
        \begin{teachernote}
        Subtract eq1 from eq2: $3a + b = 6$. Subtract eq2 from eq3: $5a + b = 0$.
        Subtract: $2a = -6 \Rightarrow a = -3$. Then $b = 0 - 5(-3) = 15$. Then $c = -1 - (-3) - 15 = -13$.
        \end{teachernote}

        $a = $ \ans{$-3$} \quad $b = $ \ans{$15$} \quad $c = $ \ans{$-13$}

  \item Write the complete quadratic model.

        \vspace{0.06in}
        $P(x) = $ \ans{$-3x^2 + 15x - 13$}

  \item State one assumption this model makes that may not hold in reality.

        \vspace{0.06in}
        \ansline{The model assumes profit follows a perfect quadratic pattern; real-world profit may fluctuate irregularly or not be symmetric.}

  \item Predict the profit when $x = 4$ (400 units). Is this prediction inside or outside
        the range of the data?

        \vspace{0.06in}
        $P(4) = -3(16) + 15(4) - 13 = -48 + 60 - 13 = $ \ans{$-1$ thousand dollars}. This is \ans{outside (extrapolation beyond)} the data range.
\end{enumerate}
\end{tcolorbox}

\end{document}
